\documentclass[11pt,letterpaper]{article}
\usepackage[margin=1in]{geometry}
\usepackage{amsmath,amssymb}
\usepackage{fancyhdr}
\usepackage{xcolor}
\usepackage{enumitem}
\usepackage[framemethod=tikz]{mdframed}
\usepackage[hidelinks]{hyperref}
\pagestyle{fancy}
\fancyhf{}
\lhead{\small AntsPhysicsLessons.com}
\rhead{\small Electricity \& Magnetism}
\cfoot{\small\thepage}
\renewcommand{\headrulewidth}{0.4pt}
\setlength{\parindent}{0pt}
\definecolor{keybg}{RGB}{240,244,250}
\newmdenv[backgroundcolor=keybg,linewidth=0pt,innerleftmargin=8pt,innerrightmargin=8pt,innertopmargin=6pt,innerbottommargin=6pt,skipabove=6pt]{answer}
\begin{document}
{\small\textsc{Unit 2 -- Gauss's Law}}\\[2pt]{\Large\bfseries Unit 2 Review}\\[8pt]Name: \rule{2.4in}{0.4pt} \hfill Date: \rule{1.4in}{0.4pt}\\[4pt]\hrule\vspace{10pt}{\small\itshape Placeholder worksheet for the site prototype. Free for classroom use.}\vspace{6pt}
\begin{enumerate}[leftmargin=*,itemsep=10pt]
\item A thin rod of length $L$ carries total charge $Q$ spread uniformly along it. Find the electric field at a point on the rod's axis a distance $a$ from the nearer end.
\vspace{2.2in}
\item A closed surface encloses point charges $q_1 = +2.0$~nC and $q_2 = -5.0$~nC (and a third charge of $+7.0$~nC sits outside). Find the net electric flux through the surface.
\vspace{2.2in}
\item An infinitely long line of charge has linear density $\lambda = 2.0\,\mu\text{C/m}$. Find $E$ at a perpendicular distance of $0.050$~m.
\vspace{2.2in}
\end{enumerate}
\end{document}